% Part: sets-functions-relations
% Chapter: arithmetization
% Section: checking-details
%
\documentclass[../../../include/open-logic-section]{subfiles}


\begin{document}

\olfileid[pt]{sfr}{arith}{check}
\olsection{Anéis e Corpos Ordenados}

Ao longo deste capítulo, afirmamos que certas definições se comportam ``como
devem''. Neste apêndice técnico, explicitaremos o que queremos dizer com isso e
(esboçaremos como) mostrar que as definições se comportam ``corretamente''.

Na \olref[int]{sec}, definimos adição e multiplicação em $\Int$. Queremos mostrar
que, assim definidas, conferem a $\Int$ a estrutura que ``quereríamos'' que
tivesse. Em particular, a estrutura em causa é a de um anel comutativo.

\begin{defn}
	Um \emph{anel comutativo} é um conjunto $S$, equipado com elementos
	distintos $0$ e $1$ e operações $+$ e $\times$, que satisfazem estas oito
	condições:
	\begin{align*}
		\emph{Associatividade}&&a + (b+ c) & = (a + b) + c \\
		&& (a \times b) \times c & = a \times (b\times c)\\
		\emph{Comutatividade}&&a + b &= b+ a  \\
		&&  a \times b&= b\times a\\ 
		\emph{Neutros}&&a + 0 &= a \\
		&& a \times 1 &= a\\
		\emph{Simétrico aditivo}&&(\exists b\in S)0&=a + b\\
		\emph{Distributividade}&&a \times (b+ c ) &= (a \times b) + (a \times c)
	\end{align*}
	Implicitamente, todas estas fórmulas estão quantificadas universalmente sobre
	$S$. Note-se que os elementos $0$ e~$1$ aqui não têm de ser os números naturais
	com o mesmo nome.
\end{defn}

Logo, para verificar que os inteiros formam um anel comutativo, basta verificar
estas oito condições. Nenhuma é \emph{difícil} de estabelecer, mas o trabalho é um
pouco laborioso. Por exemplo, eis como provar \emph{associatividade}, no caso da
adição:

\begin{proof} 
Fixemos $i, j, k \in \Int$. Então existem $a_1, b_1, a_2, b_2, a_3, b_3 \in
\Nat$ tais que $i = \equivrep{a_1, b_1}{}$ e $j =
\equivrep{a_2,b_2}{}$ e $k = \equivrep{a_3, b_3}{}$. (Por legibilidade, escrevemos
``$\equivrep{x, y}{}$'' em vez de ``$\equivrep{x, y}{\Intequiv}$''; faremos o
mesmo ao longo desta secção.) Então:
\begin{align*}
	i + (j + k) &= \equivrep{a_1, b_1}{}+(\equivrep{a_2, b_2}{} + \equivrep{a_3, b_3}{}) \\
	&= \equivrep{a_1,  b_1}{} + \equivrep{a_2+a_3, b_2+b_3}{}\\
	&= \equivrep{a_1 + (a_2 + a_3), b_1 + (b_2 + b_3)}{}\\
	&= \equivrep{(a_1 + a_2) + a_3, (b_1 + b_2) + b_3}{}\\
	&= \equivrep{a_1 + a_2, b_1 + b_2}{} + \equivrep{a_3, b_3}{}\\
	&= (\equivrep{a_1, b_1}{} + \equivrep{a_2, b_2}{}) + \equivrep{a_3, b_3}{}\\
	&= (i+j) + k
\end{align*}
apoiando-nos livremente no comportamento da adição em $\Nat$.
\end{proof}

Do mesmo modo, eis uma prova do \emph{simétrico aditivo}:

\begin{proof}
Fixemos $i \in \Int$, de modo que $i = \equivrep{a,b}{}$ para alguns $a,b \in
\Nat$. Seja $j = \equivrep{b,a}{} \in \Int$. Apoiando-nos no comportamento dos
naturais, $(a+b) + 0 = 0 + (a+b)$, logo $\tuple{a+b, b+a} \sim_\Int \tuple{0,0}$
por definição, e portanto $\equivrep{a+b, b+a}{} = \equivrep{0, 0}{} = 0_\Int$.
Assim $i + j =
\equivrep{a,b}{}+\equivrep{b,a}{}=\equivrep{a+b, b+a}{}= \equivrep{0,
0}{} = 0_\Int$.
\end{proof}

E eis uma prova da \emph{distributividade}:

\begin{proof}
Como acima, fixemos $i = \equivrep{a_1, b_1}{}$ e $j =
\equivrep{a_2,b_2}{}$ e $k = \equivrep{a_3, b_3}{}$. Então:
\begin{align*}
	i \times (j + k) 
	&= \equivrep{a_1, b_1}{} \times (\equivrep{a_2,b_2}{} + \equivrep{a_3, b_3}{})\\
	&= \equivrep{a_1, b_1}{} \times \equivrep{a_2 + a_3,b_2+b_3}{}\\
	&= \equivrep{a_1  (a_2 + a_3) + b_1  (b_2+b_3), a_1  (b_2 + b_3) + b_1 (a_2 + a_3)}{}\\
	&= \equivrep{a_1 a_2 + a_1a_3 + b_1 b_2+b_1b_3, a_1 b_2 + a_1b_3 + a_2b_1 + a_3b_1}{}\\		
%		&= \equivrep{a_1 a_2 + b_1b_2 + a_1a_3 + b_1 b_2+b_1b_3, a_1 b_2 + a_1b_3 + a_2b_1 + a_3b_1}{}\\		
	&= \equivrep{a_1a_2 + b_1b_2, a_1b_2 + a_2b_1}{} + \equivrep{a_1a_3 + b_1b_3, a_1b_3 + a_3b_1}{}\\
	&= (\equivrep{a_1, b_1}{} \times \equivrep{a_2,b_2}{}) + (\equivrep{a_1, b_1}{} \times  \equivrep{a_3, b_3}{})\\
	&= (i \times j) + (i \times k)
\end{align*}
\end{proof}

Deixamos como exercício provar as cinco condições restantes. Tendo feito isso,
mostramos que $\Int$ constitui um anel comutativo, isto é, que a adição e a
multiplicação (assim definidas) se comportam como devem.

\begin{prob}
Mostre que $\Int$ é um anel comutativo.
\end{prob}

Mas a tarefa não termina. Para além de definir adição e multiplicação em $\Int$,
definimos uma relação de ordem, $\leq$, e há que verificar que esta se comporta
como deve. Com mais pormenor, há que mostrar que $\Int$ constitui um \emph{anel
ordenado}.\footnote{Como na \olref[sfr][rel][ord]{def:linearorder}, uma
	ordem total é uma relação reflexiva, transitiva, antissimétrica e conexa.
	No contexto de relações de ordem, a conexividade é por vezes chamada
	\emph{tricotomia}, pois, para quaisquer $a$ e $b$, temos $a \leq b \lor a
	= b \lor a \geq b$.} 

\begin{defn}
Um \emph{anel ordenado} é um anel comutativo que está também equipado com uma
relação de ordem total, $\leq$, tal que:
\begin{align*}
	a \leq b &\lif a + c \leq b + c\\
	(a \leq b \land 0 \leq c) &\lif a \times c \leq b \times c
\end{align*}
\end{defn}

\begin{prob}
Mostre que $\Int$ é um anel ordenado. 
\end{prob}

Como antes, é laborioso mas rotineiro mostrar que~$\Int$, tal como construído, é
um anel ordenado. Deixamo-lo ao leitor.

Isto trata dos inteiros. Mas agora precisamos de mostrar coisas muito parecidas
para os racionais. Em particular, há que mostrar que os racionais formam um
\emph{corpo ordenado}, com as nossas definições de $+$, $\times$ e $\leq$:
\begin{defn}\ollabel{orderedfield}
Um \emph{corpo ordenado} é um anel ordenado que também satisfaz:
\begin{align*}
	\emph{Inverso multiplicativo}& & (\forall a \in S \setminus \{0\})(\exists b \in S) a\times b& = 1
\end{align*}
\end{defn}

Uma vez mostrado que $\Int$ constitui um anel ordenado, é fácil mas laborioso
mostrar que $\Rat$ constitui um corpo ordenado.

\begin{prob}
Mostre que $\Rat$ é um corpo ordenado.
\end{prob}

Tendo tratado dos inteiros e dos racionais, só restam os reais. Em particular,
há que mostrar que $\Real$ constitui um corpo ordenado \emph{completo}, isto é, um
corpo ordenado com a propriedade de completude. Agora,
\olref[cuts]{realcompleteness} estabeleceu que $\Real$ tem a propriedade de
completude. Contudo, cumpre ainda percorrer a tediosa verificação de que $\Real$ é
um corpo ordenado.

Antes de nos lançarmos a \emph{esse} exercício laborioso, há que confirmar algumas
coisas mais ``imediatas''. Por exemplo, precisamos de garantia de que
$\alpha + \beta$, tal como definida, é de fato um \emph{corte}, para quaisquer
cortes $\alpha$ e $\beta$. Segue-se uma prova desse fato:

\begin{proof}	
Como $\alpha$ e $\beta$ são ambos cortes, $\alpha + \beta = \Setabs{p
+ q}{p \in \alpha \land q \in \beta}$ é um subconjunto próprio não vazio de
$\Rat$. Suponhamos agora $x < p + q$ para alguns $p \in \alpha$ e $q \in \beta$.
Então $x - p < q$, logo $x - p \in \beta$, e $x = p + (x - p) \in
\alpha + \beta$. Logo $\alpha + \beta$ é um segmento inicial de $\Rat$.
Finalmente, para qualquer $p + q \in \alpha + \beta$, como $\alpha$ e
$\beta$ são cortes, existem $p_1 \in \alpha$ e $q_1 \in \beta$
tais que $p < p_1$ e $q < q_1$; logo $p + q < p_1 + q_1 \in \alpha +
\beta$; logo $\alpha + \beta$ não tem máximo. 
\end{proof}

Esforços semelhantes permitem verificar que $\alpha - \beta$ e
$\alpha \times \beta$ e $\alpha \div \beta$ são cortes (no último
caso, ignorando o caso em que $\beta$ é o corte zero). De novo, deixamo-lo ao
leitor.

\begin{prob}
Mostre que $\Real$ é um corpo ordenado.
\end{prob}

Mas eis um pequeno fio solto a atar. Em
\olref[cuts]{sec}, sugerimos que podemos tomar $\sqrt{2} =
\Setabs{p \in \Rat}{p < 0 \text{ ou }p^2 < 2}$. Mas há que mostrar
que este conjunto é de fato um \emph{corte}. Segue-se uma prova desse fato:

\begin{proof}
Claramente isto é um segmento inicial próprio não vazio dos racionais; basta pois
mostrar que não tem máximo. Em particular, basta mostrar que, sendo $p$ um
racional positivo com $p^2 < 2$ e $q =
\frac{2p+2}{p+2}$, se tem $p < q$ e $q^2 < 2$. Para ver que $p < q$,
basta notar:
\begin{align*}
	p^2 &< 2\\
	p^2 + 2p &< 2 + 2p\\
	p(p + 2) &< 2 + 2p\\
	p &< \tfrac{2+2p}{p+2} = q
\end{align*}
Para ver que $q^2 < 2$, basta notar:
\begin{align*}
	p^2 &< 2\\
	2p^2 + 4p + 2 &< p^2 + 4p+ 4\\
	4p^2 + 8p + 4 &< 2(p^2 + 4p + 4)\\
	(2p+2)^2 & <2(p+2)^2\\
	\tfrac{(2p+2)^2}{(p+2)^2} &< 2\\
	q^2 &< 2
\end{align*}
\end{proof}
\end{document}
